\chapter{相關研究}
\label{ch:relatedwork-zh}

本章從四個面向回顧先前研究，這些面向共同促成 RA-GARK 的設計。第~\ref{sec:rw-foundations-zh}~節介紹我們區域與全域視角所立足的兩項基礎：作為協同過濾骨幹的 LightGCN，與作為深度 KG--CF 融合典範的 KGAT。第~\ref{sec:rw-contrastive-zh}~節回顧那些在不明確閘控下對齊協同與 KG 視角的對比式 KG 方法。第~\ref{sec:rw-kgrec-zh}~節討論與本工作最為相關的合理化感知方法 KGRec，第~\ref{sec:rw-paradigm-zh}~節則以「KG 可信度」為視角，對比我們與 KGRec 的合理化典範。第~\ref{sec:rw-gating-zh}~與~\ref{sec:rw-gating-gap-zh}~節廣泛回顧深度學習中的閘控機制，並觀察到既有 KG-aware 方法皆未採用融合閘以提供架構層級的優雅退化。第~\ref{sec:rw-summary-zh}~節總結 RA-GARK 相對於上述工作的定位。

\section{基礎：LightGCN 與 KGAT}
\label{sec:rw-foundations-zh}

\textbf{LightGCN}~\cite{he2020lightgcn} 是本研究區域視角的直接前身。建立於「NGCF~\cite{wang2019ngcf} 中的非線性特徵變換對推薦準確率貢獻甚微」這項觀察上，LightGCN 將圖卷積網路精簡至最小骨幹：反覆在使用者—物品二部圖上進行正規化鄰居聚合，再將傳播後的嵌入做簡單的層平均。此一簡化在標準推薦基準上一致地優於 NGCF，且訓練更快、更穩定。其後諸多自監督方法，如 SGL~\cite{wu2021sgl} 與 DCCF~\cite{ren2023dccf}，皆在 LightGCN 式骨幹之上疊加額外的對比目標。

在 RA-GARK 中，我們\emph{原封不動}採用 LightGCN（$K=2$）作為區域視角：此分支內不施加任何 KG 相關運算。動機有二。其一，在我們的稀疏評論型 KG 上，純 LightGCN 已超越所有我們所評估的 KG-aware 基線。任何偏離此強預設的設計，皆須可被證實能改善而非劣化結果，這是本論文反覆訴諸的硬性設計約束。其二，將 KG 訊號隔離至獨立的全域視角（第~\ref{ch:methodology-zh}~章），可避免 KG 雜訊污染協同表示的梯度流。

\textbf{KGAT}~\cite{wang2019kgat} 代表 KG-aware 推薦中的典範\emph{深度融合}做法。KGAT 將使用者—物品互動圖與知識圖譜合併為一個協同知識圖（Collaborative Knowledge Graph, CKG），並以雙交互聚合器於此統一結構上傳播訊息。KG 實體嵌入因而直接參與使用者與物品表示的形成。當 KG 稠密且品質良好時，此設計相對於 CF 基線可取得顯著提升。

此處的隱含假設是 KG 可被信任，能在傳播路徑所及之處皆注入有用訊號。在我們的稀疏情境下此假設崩解：KGAT 低於純 LightGCN。其後續工作~\cite{yang2022kgcl, zou2022mcclk, yang2023kgrec} 沿用 KGAT 式聚合器，因此也繼承相同的脆弱性。RA-GARK 反轉此一設計選擇：我們完全不將 KG 與使用者—物品圖合併，而是將 KG 訊號導入一條專用旁路通道，當 KG 被證實不可靠時模型可衰減甚至全然脫接之。

\section{對比式知識圖譜方法}
\label{sec:rw-contrastive-zh}

\textbf{KGCL}~\cite{yang2022kgcl} 對 KG 引入關係層級的結構擾動，並以 InfoNCE~\cite{oord2018infonce} 對比目標對齊擾動版本與原始版本。直覺是：具資訊量的 KG 結構應對隨機擾動具有不變性，而雜訊邊則不應。KGCL 在稠密 KG 情境下顯著優於 KGAT。然而在我們的稀疏情境中，擾動後的 KG 保留的邊過少，對比目標無法提供可靠監督。

\textbf{MCCLK}~\cite{zou2022mcclk} 將對比對齊擴展至三個視角（協同、語意與結構）並對三者兩兩施加對比損失。雖然此多視角設計在稠密 KG 基準上取得 SOTA，但同樣繼承稠密 KG 的假設：在稀疏 KG 下，三個視角中有兩個（語意與結構）訊號不足，對比對齊淪為由雜訊主導而非由結構主導的正則化。

RA-GARK 同樣納入跨視角對比學習，但其角色刻意受限。我們的對比損失（$\mathcal{L}_{\mathrm{aCL}}$ 與 $\mathcal{L}_{\mathrm{uCL}}$，第~\ref{sec:cl-zh}~節）僅作為小權重（$\lambda = 0.005$）的輔助幾何正則化項，並關鍵地在 KG 側施加 stop-gradient，使對比梯度從不反向傳遞至 SVD 初始化的 aspect 嵌入。相較於 KGCL/MCCLK，此處的對比學習是表層的對齊層，而非主要的融合機制；協同與 KG 訊號的實際整合，是由一個明確的閘所執掌（第~\ref{sec:rw-gating-zh}~節）。

\section{合理化感知推薦：KGRec}
\label{sec:rw-kgrec-zh}

\textbf{KGRec}~\cite{yang2023kgrec} 是與本工作最直接相關的研究。它是第一個將合理化選擇明確化的 KG-aware 推薦器：KGRec 不將每條 KG 三元組視為同等具資訊量，而是對每條 KG 邊計算一個基於注意力的重要性分數，並以 Bernoulli dropout 隨機移除高注意力邊，迫使模型不依賴少數主導邊；其穩健性接著由原始圖與被丟邊版本之間的對比目標來強化。整套合理化機制建立於 KGAT 式聚合器之上，因此合理化作用於使用者—物品—實體統一圖的訊息傳遞路徑之內。在我們的稀疏基準上，KGRec 是四個 KG-aware 基線中最強者，但仍低於純 LightGCN。

RA-GARK 在精神上採納合理化感知典範，但在結構上多處與 KGRec 有別，下節詳述之。

\section{合理化典範：共同前提與分歧的可信度假設}
\label{sec:rw-paradigm-zh}

KGRec~\cite{yang2023kgrec} 與 RA-GARK 都認同「並非所有 KG 邊都對推薦同等貢獻、需要明確的選擇機制」這項觀察。然而二者對 KG\emph{整體可信度}所作的假設有所不同，此一假設並蔓延至各自分歧的設計選擇。

KGRec 的隱含假設是：具資訊量的邊\emph{存在}於 KG 中，因此合理化任務簡化為將其找出；隨機 dropout 配合對比學習已足以辨識並放大有用子集。KG 整體作為證據體系被假定為可靠的，僅是個別邊的權重需被調整。

RA-GARK 的前提較弱。我們不假設 KG 整體是可靠的。aspect-saliency softmax（第~\ref{sec:softmax-zh}~節）在 KG 具資訊量時選取合理化依據，但具局部偏置的融合閘（第~\ref{sec:gate-zh}~節）可在閘的梯度顯示無益時，另外地\emph{脫接整條 KG 通道}。此可信度差異展現於四項設計分歧上：

\begin{itemize}
  \item \textbf{合理化粒度。} KGRec 在邊（KG 三元組）層級進行選擇。RA-GARK 將每件物品的 KG 語意聚合為 $A=4$ 個\emph{潛在}面向槽（槽位內容由 IDF 加權商品—面向共現矩陣之截斷 SVD 投影得到）並在槽位層級進行選擇。此槽位數為模型設計參數，與每件物品所擁有的原始面向標籤數無關；粒度雖較邊層級選擇為粗，但直接產生可解讀的每物品 aspect 偏好。
  \item \textbf{選擇機制。} KGRec 採用離散 Bernoulli dropout 配合對比學習。RA-GARK 採用可微分的 softmax 注意力，能端到端訓練，並將推論期的權重作為可解釋性產物外露。
  \item \textbf{作用位置。} KGRec 在 KGAT 訊息傳遞路徑\emph{之內}施加合理化；合理化與聚合糾纏在一起。RA-GARK 將合理化置於與 LightGCN 區域視角分離的旁路通道，僅於後期由閘將兩者結合。
  \item \textbf{對 KG 可信度的處理。} KGRec 沒有任何架構機制可脫接 KG。RA-GARK 的偏置初始化融合閘正提供此一機制：訓練起始時 $\alpha^{(0)} \approx 0.993$，模型實質上即為純 LightGCN；唯當學習顯示 KG 通道能降低損失時，閘才會被打開。
\end{itemize}

表~\ref{tab:rw-kgrec-vs-ragark-zh} 將上述四項軸線並列對照。

\begin{table}[!htbp]
\centering
\caption{KGRec 與 RA-GARK 之合理化典範比較。}
\label{tab:rw-kgrec-vs-ragark-zh}
\begin{tabular}{l l l}
\toprule
軸線 & KGRec & RA-GARK \\
\midrule
合理化粒度       & KG 三元組（邊）         & $A$ 個潛在面向槽 \\
選擇機制         & Bernoulli dropout + CL  & Softmax 注意力 \\
作用位置         & KGAT 訊息傳遞之內       & 旁路通道，後期融合 \\
KG 可信度處理    & 無脫接機制              & 偏置初始化融合閘 \\
\bottomrule
\end{tabular}
\end{table}
\FloatBarrier

因此我們不將 RA-GARK 定位為對 KGRec 的工程性改良，而是基於對 KG 可信度不同前提之另一種設計。此一區別在經驗上是否重要，則於第~\ref{ch:experiments-zh}~章中作答。

\section{深度學習中的閘控機制}
\label{sec:rw-gating-zh}

RA-GARK 的融合閘有兩個 KG-aware 推薦之外的近親，二者皆啟發本研究設計。

\textbf{Highway Networks}~\cite{srivastava2015highway} 引入可學習的閘 $T(\mathbf{x})$，於變換訊號 $H(\mathbf{x})$ 與恆等跳接間進行插值：$\mathbf{y} = T(\mathbf{x}) \cdot H(\mathbf{x}) + (1 - T(\mathbf{x})) \cdot \mathbf{x}$。其關鍵設計選擇為對 $T$ 的負偏置初始化：訓練之初 $T \approx 0$，網路表現為恆等映射。其後訓練僅在變換有益處之處才打開閘門。此設計原本是為了應對極深網路訓練的困難，但「將閘偏置至已知安全的預設，僅在學習認為合理時才覆寫」此結構性樣式，遠遠超出深度的脈絡而具有普適性。RA-GARK 將此樣式由網路內跳接移植至雙視角融合：融合閘上的 $+5$ 偏置扮演與 Highway 之負偏置相同的角色，對 KG 通道提供優雅退化保證。

\textbf{專家混合（Mixture-of-Experts）閘控}，以多任務推薦中的 MMoE~\cite{ma2018mmoe} 與 PLE~\cite{tang2020ple} 為代表，以閘控網路對多個專家塔加權。雖然概念上類似多管線選擇，但兩項重要性質與 RA-GARK 不同。其一，MMoE/PLE 中的專家為\emph{同質候選者}（皆為相同架構、僅以任務特定目標訓練的實例）因此並無「某專家本質上比另一者較不可靠」的概念。其二，這些架構中的閘採對稱初始化（訓練起始不偏向任何特定專家），適用於對稱專家情境，但不適用於我們不對稱的情境，在我們的設定中，一條通道（LightGCN）已知是強預設，另一條（KG）已知有時並不具資訊量。

另有一條工作線是對比式雙視角推薦器，例如 SGL~\cite{wu2021sgl} 與 DCCF~\cite{ren2023dccf}，這類方法以擾動單一圖（邊 dropout、節點 dropout、隨機漫步）來建構多個視角。此類方法中兩個視角位於\emph{相同訊號空間}，對比對齊已足以整合二者。RA-GARK 的兩個視角則在結構上異質：區域視角為使用者—物品二部圖，全域視角為物品—aspect 的 KG。僅靠對比學習進行隱式對齊，無法調整兩個異質訊號空間的相對權重；明確的閘是必要的。

\section{缺口：KG-Aware 推薦中沒有融合閘}
\label{sec:rw-gating-gap-zh}

儘管閘控在深度學習中已相當成熟，我們所調研的 KG-aware 推薦文獻中，未見任何方法將融合閘作為協同訊號與 KG 訊號之間的整合機制。表~\ref{tab:rw-gating-gap-zh} 摘要了四個被基準測試的基線與另外三個知名 KG-aware 方法（CKAN~\cite{wang2020ckan}、MKR~\cite{wang2019mkr}、KGIN~\cite{wang2021kgin}）所採用的整合機制。在這個集合中，KG 訊號是透過直接實體傳播、對比對齊、邊層級 dropout、注意力、跨特徵單元或意圖解耦模組來整合。沒有任何機制允許 KG 通道在訊號不可靠時被\emph{結構性地脫接}。

\begin{table}[!htbp]
\centering
\caption{代表性推薦方法中的 KG 整合機制。}
\label{tab:rw-gating-gap-zh}
\begin{tabular}{l l c}
\toprule
方法 & KG 整合機制 & 可脫接？ \\
\midrule
KGAT~\cite{wang2019kgat}      & 直接實體訊息傳播       & 否 \\
KGCL~\cite{yang2022kgcl}      & 擾動 KG 對比學習       & 否 \\
MCCLK~\cite{zou2022mcclk}     & 多視角對比對齊         & 否 \\
KGRec~\cite{yang2023kgrec}    & 邊層級 dropout + KGAT 聚合器 & 否 \\
CKAN~\cite{wang2020ckan}      & 協同—知識注意力        & 否 \\
MKR~\cite{wang2019mkr}        & 跨特徵轉移單元         & 否 \\
KGIN~\cite{wang2021kgin}      & KG 上的意圖解耦        & 否 \\
\midrule
\textbf{RA-GARK}（本研究）    & \textbf{偏置初始化的融合閘} & \textbf{是} \\
\bottomrule
\end{tabular}
\end{table}
\FloatBarrier

我們刻意作出一個範圍狹義的主張。我們不主張 RA-GARK 是 KG-aware 推薦中首個使用「任何閘」的方法，因 KG-aware 文獻浩繁，無法窮舉列舉。可辯護的主張為：\emph{在我們所檢視的方法中，沒有任何一個以偏置初始化的融合閘於稀疏或不可靠 KG 下提供架構優雅退化}。我們在第~\ref{ch:experiments-zh}~章的消融顯示，此一初始化在本設定中確實是不可忽略的設計選擇。

\section{注意力正規化}
\label{sec:rw-attention-norm-zh}

一個橫跨多領域的相關問題是：在選擇合理化依據時，應對注意力權重套用 sigmoid 或 softmax？文獻並未匯聚於單一答案。DIN~\cite{zhou2018din} 對使用者行為歷史套用 sigmoid 注意力，所基於的假設是每筆歷史行為與候選物品的相關性彼此獨立。NAIS~\cite{he2018nais} 對歷史互動套用 softmax 注意力以衡量物品相似度，AFM~\cite{xiao2017afm} 則對特徵交互套用 softmax。各工作中的選擇皆以其注意力所依附的語意假設為依據。

我們在第~\ref{ch:experiments-zh}~章報告一項單一資料集的觀察：對於我們的 aspect-saliency 合理化，將 softmax 替換為 sigmoid 會產生我們所量測到最大的單一元件變化。我們將其解讀為支持「正規化應與合理化的語意假設相符」此一設計啟發，但明確將此觀察定位為待多資料集複製驗證的假設（第~\ref{ch:conclusion-zh}~章）。

\section{總結與 RA-GARK 的定位}
\label{sec:rw-summary-zh}

表~\ref{tab:rw-summary-zh} 從四個設計面向匯整比較。在本章所討論的方法中，RA-GARK 是唯一同時將純 LightGCN 區域聚合器與閘控後期融合的 KG 訊號相結合者，也是唯一將合理化作用於物品—aspect 粒度而非 KG 邊粒度者。

\begin{table}[!htbp]
\centering
\caption{KG-aware 推薦方法之設計比較。}
\label{tab:rw-summary-zh}
\begin{tabular}{l l l l}
\toprule
方法 & 區域聚合器 & KG 整合 & 合理化 \\
\midrule
KGAT       & 雙交互           & 直接傳播             & 無 \\
KGCL       & KGAT 式          & 傳播 + CL            & 無 \\
MCCLK      & 多視角           & 多視角 CL            & 無 \\
KGRec      & KGAT 式          & 傳播 + 邊 dropout    & 邊 Bernoulli \\
\midrule
\textbf{RA-GARK} & \textbf{純 LightGCN} & \textbf{閘控後期融合} & \textbf{Aspect softmax} \\
\bottomrule
\end{tabular}
\end{table}
\FloatBarrier

我們並不將 RA-GARK 定位為相對於四個基線的嚴格改進。RA-GARK 在設計空間中佔據一個不同的位置，其特徵在於將 KG 視為可被明確閘控的旁路通道，而非必要的管線組件。此選擇在我們的稀疏 KG 情境中產出強勁結果，而基於閘控的優雅退化性質，在 KG 品質不均或模型部署於 KG 密度差異各異的多個領域時，或許亦能派上用場。然而，當 KG 稠密且品質一致地可靠時，KGAT 或 KGRec 這類深度融合方法可能仍較佳，因為後期融合的設計犧牲了深度融合架構所具有的部分表徵表達力。
